\subsection{逆tangent正Pfaffianの共鳴核は全て一次因子を持つ}
\begin{theorem}\label{thm:h3-order2-22-reverse-resonant-empty}
逆tangent・非退化正top・正Pfaffianのatlas
\[
 (z_0,z_1,z_2,z_3,z_4,z_5)
 =(z_0,z_1,-z_4/2,z_4/2,z_4,z_5),\quad
 z_0z_5-z_1z_4-z_4^2/4=1
\]
のうち、共鳴 $z_4=1$ の全領域にH3は存在しない。
この領域では正常三次核は二次元であるが、全切断が一次因子を持つためである。
$z_4\ne1$ の全被覆・正常核や、平方正topのrank2 cross境界はこの定理に含めない。
\end{theorem}
\begin{proof}
$z_5\ne0$ なら正側の保存torusで $z_5=1$ とでき、全包絡は
\[
 (z_0,z_1,z_2,z_3,z_4,z_5)
 =\bigl(\beta+5/4,\beta,-1/2,1/2,1,1\bigr)
\]
と書ける。$D=(\beta-1)(\beta+1)(2\beta+1)(4\beta+5)$ と置く。
元の全三次・四次関係から商線型汎関数を作ると、その全係数は $k[\beta,D^{-1}]$ に属する。
元の全関係を消す恒等式を検証し、自由語上に恒等部分行列を持つので、
AS4候補点ではHilbert値により正三次商の座標を正確に与える。
以下の座標は順に
\[
 (x^3,y^3,u^2x,u^2y,uvx,uvy,vux,vuy,v^2x,v^2y)
\]
である。$d=\beta+5/4$ とすると、正常核の二列は
\[
 W_0=\begin{pmatrix}-1\\-d^{-1}\\\beta+2\\1\\-1\\1\\1\\-1\\1\\0\end{pmatrix},\qquad
 W_1=\begin{pmatrix}-d\\-d^{-1}\\d\\\beta+2\\-d\\1\\d\\-1\\0\\1\end{pmatrix}.
\]
四次商汎関数で作った正常行列について、二つの $8\times8$ minorのgcdは定数倍の $D^2$。
従って全 $D\ne0$ で核は正確にこの二次元となり、例外値の未処理部分を一般点計算へ隠していない。

この二次元核を零と主張する代わりに、全三次に一次因子があることを示す。
正二次商の固定基底を
\[
 (x^2,xy,yx,y^2,uv,v^2)
\]
とする。偶二次関係は $u^2,vu$ を単位係数で消せるため、全係数でこの六語が基底となる。
入力JSONには、$\theta$ に関して二次、$\beta$ に関して多項式の六係数 $Q_2(\beta,\theta)$ が与えられる。
元の関係から得た三次商座標で
\[
 (x+\theta y)Q_2(\beta,\theta)
 =g_0(\beta,\theta)W_0+g_1(\beta,\theta)W_1
\]
を厳密検証する。右辺の二係数は $\theta$ 三次以下であり、$\beta=0$ で比例しない。
従って一般の $\beta$ では $\mathbb P^1_{\theta}$ から正常pencilへの写像は非定数で、像は稠密である。

この一般点の因数分解を全 $D\ne0$ に延ばす理由を明記する。
商座標の乗法行列 $L(\ell)$ に対し、
\[
 [L(\ell)q]\wedge[c_0W_0+c_1W_1]=0
\]
という比例条件を
$D\ne0$ 上の $\mathbb P^1_{\ell}\times\mathbb P^5_q\times\mathbb P^1_c$ に課す。
各成分は多重斉次な多項式方程式で、$\mathbb P^1_c$ への射影はproperである。
一般fiberで稠密な像を持つため、その閉像は既約な全
$(D\ne0)\times\mathbb P^1_c$ に等しい。
この閉incidenceは零積も形式上許すが、仮に対象がH3なら包絡はdomainである。
従って非零 $\ell$ と非零正二次式 $q$ の積は非零であり、
上の比例条件は実際に $c_0W_0+c_1W_1$ の一次因子を与える。
商座標と実際の三次商の同定には、当該H3候補のHilbert値だけを使用した。
全係数でのAS4存在やintrinsic族降下を先に仮定する議論ではない。

$D=0$ の四値 $1,-1,-1/2,-5/4$ は別に元の関係から消去する。
各点の正常行列の階数は8、同じ $Q_2$ による三次像の係数spanは二次元である。
各点での非定数因子incidenceとpropernessにより、全正常pencilがやはり一次因子を持つ。
これで $z_5\ne0$ の全点を扱った。

$x,y$ の交換と $v\mapsto-v$ はこの共鳴atlasで $z_0,z_5$ を交換する。
従って $z_5=0,z_0\ne0$ も上に含まれる。
最後の $z_0=z_5=0,z_1=-5/4$ では、正側が負側を互いに逆な
\[
 f=\begin{pmatrix}1/2&3/4\\-1&1/2\end{pmatrix},\qquad g=f^{-1}
\]
でnormalizeし、$f^3=-I$ である。
正常核は $kx^3+ky^3$ で、次の恒等式が成立する：
\[
 \ell^3-\theta C\ell=x^3+\theta^3y^3,\qquad
 \ell=x+\theta y,\quad C=-\tfrac34u^2-v^2.
\]
$C$ は中心二次式なので、左辺は $\ell(\ell^2-\theta C)$ と分解する。
$\theta\mapsto\theta^3$ は代数閉体上全射であり、$y^3$ 自身にも因子があるから、
この最後のpencilも全て分解する。

以上の各場合で、三次関係による商の一次・二次因子はどちらも非零なので、
商は非domainとなりH3に反する。
非退化正topのJordan排除は既知である。
検証器
\[
\texttt{verify\_cubic\_order2\_22\_reverse\_resonant.py}
\]
は全元関係に対する商汎関数、二minor gcd、六係数因子恒等式、四例外値、
軸交換、及び最後の三次norm恒等式を入力から厳密再生する。
独立査読や共鳴外の完全判定は主張しない。
\end{proof}
